% SPDX-License-Identifier: CC-BY-4.0
% Copyright (c) 2026 Martin Schröder <info@swedishembedded.com>

\section{Method}\label{sec:method}\label{sec:rsi-design}

\paragraph{Terms.} A \emph{series} is a registered experiment of four \emph{lineages}, each a chain
of three \emph{generations} from the base that differs from its siblings only in its seed; the
\emph{champion} is the model a generation's material comes from; a \emph{step} compares a generation with
the one before. The registrations call the series S1 to S4 (``recursive self-improvement''); we call
them S1 to S4, because what they test is multi-round self-training on self-proposed, oracle-labelled
material, and whether it is recursive in the sense that an improved model is a better improver is what
the matched-data control decides (\cref{sec:followups}).\cl{RS01}

\paragraph{The learner.} The task family is predicting what a Lua~5.1 program prints. The learner is
the loop of \cref{sec:documentation} reduced to this one kind: the policy, Qwen3-4B with thinking off,
writes programs about sections of the Lua~5.1 manual, the pinned Lua~5.1.5 interpreter runs each
program and its output becomes the reference, and the policy is trained on its own programs with the
interpreter's output as the target. From S2 on every model call and every training run goes through
Brain's API: one \code{brain serve} per lineage serves the proposals, the predictions and the dev and
retention exams, and trains each candidate as a Brain training job.\cl{RS09,RS10} A lineage is a chain
of generations from the base that differs from its siblings only in its seed; each series runs four
lineages of three generations. Inside the loop a candidate is accepted by a dev gate on 400 units
(100 generator programs and 300 programs the base proposed, held out from training) with a retention
set of 300 units; neither holds an exam item.\cl{RS09} The generator programs come from a separate seeded generator of six levels
(a sum, a product, a short string transformation, a loop, a three-item table sort, an arithmetic
expression) that shares no code or templates with the exam's generator, although two of its levels call
functions for which the exam has templates (\code{string.upper}, \code{string.rep}, the length
operator, \code{table.sort}).\cl{RS51}

\begin{figure}[tbp]
\centering
\begin{tikzpicture}[fig]
\tikzset{cell/.style={text width=28mm}}
\def\dx{40mm}
% row 0: the source
\node[data, cell] (manual) at (\dx,0) {Lua 5.1 manual,\\one section per request};
% row 1
\node[proc=okSky!22, cell] (champ) at (0,-21mm) {champion of\\generation $g-1$\\(the base at $g=1$)};
\node[proc=okSky!22, cell] (prop) at (\dx,-21mm) {champion proposes\\a program (own\\word and number)};
\node[verif, cell] (lua) at (2*\dx,-21mm) {Lua 5.1.5 runs it\\twice; the equal\\output is the label};
\node[proc=okSky!22, cell] (pred) at (3*\dx,-21mm) {champion predicts\\the output $k=4$\\times, $T=1.0$};
% row 2, right to left
\node[decide, aspect=2.2, inner sep=0.5pt, text width=15mm] (front) at (3*\dx,-49mm) {any of the\\4 wrong?};
\node[data, cell] (union) at (2*\dx,-49mm) {frontier union:\\300, 900, 2100\\programs};
\node[proc=okSky!22, cell] (lora) at (\dx,-49mm) {fresh LoRA from\\the base, 2 passes:\\75, 225, 525 steps};
\node[verif, cell] (gate) at (0,-49mm) {learner's gate:\\dev 400 units,\\retention 300};
% row 3
\node[data, cell] (hand) at (\dx,-76mm) {hand-off: every\\trained program,\\candidate record};
% flows
\draw[flow] (champ) -- (prop);
\draw[flow] (manual) -- (prop);
\draw[flow] (prop) -- (lua);
\draw[flow] (lua) -- (pred);
\draw[flow] (pred) -- (front);
\draw[flow] (front) -- node[lab, above] {yes} (union);
\draw[flow] (union) -- (lora);
\draw[flow] (lora) -- (gate);
\draw[flow] ($(gate.south)+(6mm,0)$) |- node[lab, pos=0.3, right] {accepted} (hand.west);
\draw[ctrl] (gate) -- node[lab, right, align=left] {new\\champion} (champ);
\draw[refuse] ($(gate.south)+(-6mm,0)$) -- ++(0,-9mm) node[reflabel, below, align=center] {rejected:\\champion stays};
\draw[refuse] (front.south) -- ++(0,-8mm) node[reflabel, below] {all 4 right: not trained};
\draw[refuse] (lua.north) -- ++(0,6mm) node[reflabel, above, align=center] {refused: no stable output,\\a repeat, a nondeterministic call};
\draw[ctrl] (hand.east) -- node[lab, above] {to the examiner} ++(30mm,0);
\end{tikzpicture}
\caption{One generation of S4, the series that met its criterion, read clockwise from the champion. The champion proposes programs about
the manual and the interpreter labels them; a program joins the frontier when the champion predicts it
wrong in at least one of 4 sampled attempts (300, 600 and 1200 new programs at generations 1 to 3); the
candidate is a fresh LoRA trained from the base on the union of every generation's frontier; the
learner's gate decides whether it becomes the next champion, and the examiner, not the gate, decides
the registered outcome. Dotted arrows are what each stage refuses.}\label{fig:cycle}
\end{figure}

\paragraph{The exam.} The exam is 889 Lua~5.1 programs written by code from 25 standard-library
templates over randomised and edge-case arguments, the items the base never answered in four sampled
attempts under a neutral request at selection (at S4's serving tier it answers 12 of them); the reference of each is what the interpreter prints (twice,
equal).\cl{RS03} Because it was selected as what the base could not do, a count of items right on it
measures progress on that hard subset, not a model's accuracy on Lua programs in general. It was
selected once, in S1, and every series uses the same file (its SHA-256 begins
\code{60ac3ae44651} in every verdict).\cl{RS20,RS31} No model writes or answers an exam item to build
it, and no run of the learner reads it. The reuse has a cost: S1 graded this exam item by item with
its own scripts, and each later design was chosen after the aggregate results of the series before it
on the same exam, so across the series the exam is a holdout reused adaptively, a setting in which a
reused holdout can be overfitted~\citep{holdout2015}\cl{CT22}; each registration on its own is clean.
A re-examination of S4's champions on a freshly drawn exam is registered for this
reason.\cl{RS39}

\paragraph{The examiner.} From S2 on the exam is held by an examiner, a separate program operated
beside the learner by the same operators; it is independent of the learner, not of the authors. The
learner's only contact with it is a one-way hand-off (\cref{fig:examiner}).\cl{RS04} For each lineage and generation the learner writes a directory with
the list of every program the candidate's chain trained on and, last, a candidate record (base model,
base digest, adapter path and digest, serving flags, commits). The examiner then
\begin{enumerate}
\item serves the base and every candidate itself, on one pinned Brain serving binary (SHA-256
  \code{d11e9134...}, a different Brain build from the learner's commit \code{796871647}) and, in S4, at one KV-cache tier (f32), and marks a verdict invalid when the
  served base is not the declared checkpoint (identified from the checkpoint's own files) or the
  adapter does not have the declared digest; that build's source is not in Brain's published
  history, so the reference examiner of \cref{sec:repro} serves on the learner's Brain instead;\cl{RS06,RS32}
\item excludes from each step's comparison every exam item that equals, exactly or after removing
  comments and whitespace, any program on either compared model's hand-off list (the programs its
  chain trained on and, beside them, its held-out dev and retention programs), and reports how many
  it excluded; it also reports a near-duplicate count whose similarity measure the verdicts do not
  define, and a program of the same template with other arguments is not excluded;\cl{RS05}
\item asks each item once under the direct-completion protocol: the program and its question, followed
  by the sentence ``Answer with exactly what it prints, in one fenced code block.'' (the request S1
  registered), as the only user message, no system prompt, one greedy completion with thinking off and, in S4, at most
  2048 tokens; the answer is the last fenced block of the reply, or the whole reply if it has none,
  compared with the reference token by token per line with whitespace runs collapsed and blank lines
  ignored; a reply the token limit cut is graded as written;\cl{RS32,RS36}
\item pairs the candidate with its predecessor item by item (a win: the candidate is right and the
  predecessor wrong; a loss: the reverse), computes the exact one-sided sign test over the discordant
  items, and decides the registered rule: a step holds when wins exceed losses at $p<0.05$, a lineage
  succeeds when all three of its steps hold, and the series meets its criterion when at least three of
  the four lineages succeed;\cl{RS02}
\item returns, per step, a verdict with aggregates only (items right, wins, losses, $p$, whether the
  step holds, exclusion and overlap counts, integrity flags) and, over the series, a plan verdict;
  each is signed with the examiner's HMAC-SHA256 key, so the learner, which does not hold the key,
  cannot alter or forge one; the key protects against the learner, not against the authors. Each verdict
  also carries the paired gain in the fraction of items right with a 95\% interval clustered over the
  exam's 25 templates (\cref{sec:rsi-robust}). No item text and no item-level grade leaves the
  examiner.\cl{RS04,RS31,RS38}
\end{enumerate}
The items themselves are not secret: the exam file was released with S1's records, so the
separation rests on the learner's code never opening it and on the exclusion of overlapping programs,
not on secrecy.\cl{RS03}

\begin{figure}[tbp]
\centering
\begin{tikzpicture}[fig]
\tikzset{cell/.style={text width=29mm}}
% learner column
\node[proc=okSky!22, cell] (gen) at (0,0) {generation $g$:\\propose, label,\\train via Brain};
\node[verif, cell] (lgate) at (0,-22mm) {learner's gate:\\dev 400 units,\\retention 300};
% hand-off, between the two
\node[data, cell] (hand) at (37mm,-22mm) {\textbf{hand-off}:\\trained programs,\\then candidate};
% examiner, down the first column and up the second
\node[data, cell] (exam) at (76mm,0) {\textbf{held-out exam}:\\889 items the base\\missed at selection};
\node[verif, cell] (excl) at (76mm,-22mm) {drop each item\\equal to a listed\\program (also with\\comments removed)};
\node[proc=okOrange!20, cell] (serve) at (76mm,-46mm) {serve all models:\\one pinned build\\and KV tier,\\digests checked};
\node[proc=okOrange!20, cell] (grade) at (113mm,-46mm) {direct completion:\\greedy decoding,\\thinking off,\\last fenced block};
\node[decide, aspect=2, inner sep=0.5pt, text width=14mm] (rule) at (113mm,-22mm) {$W>L$,\\$p<0.05$?};
\node[data, cell] (verd) at (113mm,0) {\textbf{signed verdict}:\\counts, $p$, holds;\\no item, no grade};
\begin{scope}[on background layer]
\node[bound, fit=(gen)(lgate), inner sep=2.5mm] (L) {};
\node[bound, fit=(exam)(serve)(grade)(verd)(rule), inner sep=2.5mm] (E) {};
\end{scope}
\node[boundlabel] at (L.north west) {Learner};
\node[boundlabel] at (E.north west) {Examiner (separate program)};
\draw[flow] (gen) -- (lgate);
\draw[flow] (lgate) -- (hand);
\draw[flow] (hand) -- (excl);
\draw[flow] (hand.south) |- (serve.west);
\draw[flow] (exam) -- (excl);
\draw[flow] (excl) -- (serve);
\draw[flow] (serve) -- (grade);
\draw[flow] (grade) -- node[lab, left, align=right] {wins,\\losses} (rule);
\draw[flow] (rule) -- node[lab, left] {step holds?} (verd);
\draw[refuse] (exam.west) -- node[reflabel, above, align=center] {no path back: no item,\\grade or verdict} (gen.east);
\end{tikzpicture}
\caption{The learner and the examiner. The only channel is one-way: per lineage and
generation the learner writes the list of every program the candidate's chain trained on (with its
held-out dev and retention programs) and then the candidate record (base, adapter and their digests).
The examiner holds the exam, drops items that match a listed program (exactly, or after removing comments and whitespace), serves every model itself
on one pinned build, grades by the direct-completion protocol, applies the registered per-step rule
and signs aggregates only, which are published with the learner's records after the series. The
dotted path does not exist.}\label{fig:examiner}
\end{figure}

\paragraph{What is compared.} Every comparison is within a lineage and between consecutive
generations, on the items no compared model trained on, so a model's gain is measured against the
model it was trained from. The registered decision is per lineage. The four lineages share the exam's
items, so a sign test pooled over lineages treats correlated pairs as independent; the pooled figures
are reported as a description, not as a test. The examiner's base-against-base null replicate, run before S4, answered
12 of the 889 items on both sides with 0 discordant items and 0 failed requests: the protocol is
deterministic at the serving build and tier, so a discordant item is a difference between the
models.\cl{RS07}


\paragraph{Three levels of claim.} The item-level sign test supports a claim about two adapters: this
adapter answers more of these items than its predecessor. The template-clustered interval supports a
claim about programs of these templates. Only replication over seeds supports a claim about the
procedure, and with four lineages that is a count, not a test: 4 of 4 has a one-sided sign-test $p$ of
0.0625, and the generation-3 champions answered 219.0 items on average with a standard deviation of 20.4
between seeds.\cl{RS53} Each result below is stated at its level.
